\let\textcircled=\pgftextcircled
\chapter{Analyse et sp\'ecification des besoins}
\label{chap:analyse}

\textit{\initial{C}e chapitre formalise ce que le syst\`eme doit faire, avant toute
consid\'eration sur la mani\`ere dont il le fera. Il identifie les acteurs, \'enonce les
exigences fonctionnelles et non fonctionnelles, et les traduit en cas
d'utilisation dont les plus critiques sont d\'ecrits en d\'etail.}

\mtcselectlanguage{french}
\minitoc
\newpage

\section{D\'emarche de mod\'elisation}
\subsection{Choix du processus}
Le processus retenu est le \textbf{2TUP} (\emph{Two Track Unified Process}),
d\'eclinaison du processus unifi\'e qui s\'epare explicitement la branche fonctionnelle
--- ce que le m\'etier demande --- de la branche technique --- ce que
l'infrastructure impose --- avant de les faire converger en conception.

Ce choix est motiv\'e par la nature du projet. Les contraintes techniques y ne sont
pas subordonn\'ees au fonctionnel : la souverainet\'e des donn\'ees et le
fonctionnement hors ligne \'ecartent d'embl\'ee des solutions qui satisferaient
parfaitement le besoin m\'etier. Traiter les deux branches en parall\`ele \'evite de
concevoir une solution fonctionnellement correcte mais non d\'eployable.

\begin{figure}[H]
\centering
\begin{diagramme}[node distance=6mm]
  \node (bf)  [action=3.4cm, fill=umltete] {Branche fonctionnelle\\[2pt]\scriptsize
        Capture des besoins m\'etier\\ Analyse\\ Cas d'utilisation};
  \node (bt)  [action=3.4cm, fill=umltete, right=42mm of bf] {Branche technique\\[2pt]\scriptsize
        Contraintes de d\'eploiement\\ Architecture\\ Choix des composants};
  \node (conc) [action=4cm, below=16mm of $(bf)!0.5!(bt)$] {Conception pr\'eliminaire\\ puis d\'etaill\'ee};
  \node (real) [action=4cm, below=8mm of conc] {R\'ealisation};
  \node (val)  [action=4cm, below=8mm of real] {Recette};
  \draw[flot] (bf.south) -- (conc.north west);
  \draw[flot] (bt.south) -- (conc.north east);
  \draw[flot] (conc) -- (real);
  \draw[flot] (real) -- (val);
  \draw[dependance] (val.east) to[out=0,in=0,looseness=1.6] node[role,right=1mm]{it\'eration} (conc.east);
\end{diagramme}
\caption{D\'emarche en Y du processus 2TUP appliqu\'ee au projet}
\label{fig:2tup}
\end{figure}

\subsection{Recueil des besoins}
% A COMPLETER : preciser le dispositif reel de recueil.
Le recueil s'est appuy\'e sur [nombre] entretiens semi-directifs conduits aupr\`es
des analystes du CIRT, sur l'observation de [nombre] traitements d'incidents
r\'eels et sur l'analyse du document interne de \emph{classification des r\'eactions
autonomes par type d'attaque}, qui constitue la r\'ef\'erence m\'etier du projet.

\section{Identification des acteurs}
Un acteur est un r\^ole tenu vis-\`a-vis du syst\`eme, non une personne. Le m\^eme agent
peut en tenir plusieurs successivement --- un administrateur est, dans
l'organisation du CIRT, un analyste promu.

\begin{table}[H]
\centering
\caption{Acteurs du syst\`eme et responsabilit\'es associ\'ees}
\label{tab:acteurs}
\small
\begin{tabular}{@{}p{2.6cm}p{2.2cm}p{8.5cm}@{}}
\toprule
\textbf{Acteur} & \textbf{Nature} & \textbf{Responsabilit\'e vis-\`a-vis du syst\`eme} \\
\midrule
Analyste & Humain, principal &
  Supervise \emph{a posteriori}. Consulte le journal et le portefeuille, d\'eclenche
  une annulation manuelle, qualifie une menace hors catalogue. Ne valide jamais
  une action avant son ex\'ecution. \\
Administrateur & Humain, principal &
  Analyste promu. Ajoute la d\'efinition de la politique de r\'eponse, la gestion du
  catalogue de r\'eversibilit\'e, le r\'eglage du coupe-circuit et la gestion des comptes. \\
Super-admin. & Humain, principal &
  Cr\'e\'e \`a l'installation. Ajoute la promotion et la r\'evocation des administrateurs. \\
D\'ecideur & Humain, secondaire &
  Lecture strat\'egique stricte : portefeuille et journal. Aucune action. \\
Source de d\'etection & Syst\`eme, primaire &
  \'Emet les observations : SIEM, sonde r\'eseau, EDR, journaux syst\`eme. \\
Ordonnanceur interne & Syst\`eme, primaire &
  D\'eclenche les traitements p\'eriodiques : sondes d'infrastructure, boucle de
  contr\^ole post-action, rejeu de la file du mode d\'egrad\'e. \\
\'Equipement cible & Syst\`eme, secondaire &
  Subit les actions : pare-feu, EDR, annuaire, r\'esolveur DNS, sauvegarde. \\
\bottomrule
\end{tabular}
\end{table}

\section{Besoins fonctionnels}
\label{sec:besoins-fonctionnels}
Les exigences sont num\'erot\'ees \texttt{EF-\textit{nn}}. Cette num\'erotation est
conserv\'ee sans changement jusqu'\`a la matrice de tra\c cabilit\'e du
chapitre~\ref{chap:conception} et jusqu'aux noms des tests de recette : c'est ce
qui permet de v\'erifier, exigence par exigence, qu'aucune n'a \'et\'e perdue en route.

\begin{table}[H]
\centering
\caption{Exigences fonctionnelles}
\label{tab:exigences-fonctionnelles}
\scriptsize
\begin{tabular}{@{}llp{9.2cm}@{}}
\toprule
\textbf{R\'ef.} & \textbf{Groupe} & \textbf{\'Enonc\'e} \\
\midrule
EF-18 & Ingestion & Normaliser les remont\'ees de sources h\'et\'erog\`enes en un sch\'ema pivot unique. \\
EF-19 & Ingestion & \'Eliminer les remont\'ees multiples d'une m\^eme observation. \\
EF-20 & Ingestion & Corr\'eler les \'ev\'enements en incidents sur une cl\'e et une fen\^etre temporelle. \\
\addlinespace
EF-03 & Analyse & Enrichir le contexte d'un incident \`a partir d'un fonds documentaire. \\
EF-04 & Analyse & \textbf{N'engager aucune action lorsque le contexte n'est pas document\'e.} \\
EF-08 & Analyse & \'Etablir un profil comportemental de r\'ef\'erence par entit\'e. \\
EF-09 & Analyse & D\'etecter un \'ecart significatif \`a ce profil. \\
EF-10 & Analyse & Rendre le score d'anomalie explicable. \\
\addlinespace
EF-05 & D\'ecision & S\'electionner l'action corrective selon la famille de menace. \\
EF-06 & D\'ecision & Param\'etrer l'action \`a partir des indicateurs de l'\'ev\'enement. \\
EF-14 & D\'ecision & \textbf{N'autoriser en autonomie que les actions r\'eversibles.} \\
EF-15 & D\'ecision & \textbf{Appliquer une politique r\'edig\'ee en langage naturel, compil\'ee a priori.} \\
EF-26 & D\'ecision & \textbf{Suspendre globalement l'autonomie au-del\`a d'un seuil d'\'echec.} \\
\addlinespace
EF-07 & Ex\'ecution & \textbf{Ex\'ecuter l'action sans validation humaine pr\'ealable.} \\
EF-13 & Ex\'ecution & \textbf{Notifier l'analyste apr\`es coup, sans bloquer l'ex\'ecution.} \\
EF-25 & Ex\'ecution & \textbf{Surveiller l'effet de l'action et l'annuler si la cible se d\'egrade.} \\
\addlinespace
EF-21 & Surveillance & Surveiller la disponibilit\'e des services supervis\'es. \\
EF-22 & Surveillance & Param\'etrer des seuils de service par cible. \\
EF-23 & Surveillance & Traiter une d\'egradation d'infrastructure comme un incident. \\
\addlinespace
EF-11 & Consultation & Consulter l'\'etat courant du syst\`eme. \\
EF-12 & Consultation & Consulter le journal des d\'ecisions. \\
EF-16 & Consultation & Prioriser le portefeuille d'incidents. \\
EF-17 & Consultation & Produire les indicateurs de pilotage. \\
\addlinespace
EF-27 & Hors catalogue & \textbf{Confiner une menace inconnue \`a partir des seuls indicateurs observ\'es.} \\
EF-28 & Hors catalogue & \textbf{Alerter de mani\`ere persistante sur un geste \`a effet durable non engag\'e.} \\
EF-29 & Hors catalogue & \textbf{Proposer une qualification soumise \`a validation humaine, et en enrichir le catalogue.} \\
\bottomrule
\end{tabular}
\end{table}

\begin{tcolorbox}[colback=umlalerte!5, colframe=umlalerte, boxrule=0.5pt, sharp corners,
                  fontupper=\small, title=\textbf{Inversion assum\'ee par rapport \`a la version pr\'ec\'edente du cahier des charges}]
Les exigences EF-07 et EF-13 inversent la position ant\'erieure, qui interdisait
toute ex\'ecution automatique et imposait la validation manuelle de chaque
recommandation. Cette inversion n'est pas un assouplissement : elle est
compens\'ee par EF-04, EF-14, EF-15, EF-25 et EF-26, qui sont autant de
\emph{refus} et de \emph{corrections} substitu\'es \`a la relecture humaine. Aucun de
ces m\'ecanismes ne r\'eintroduit d'attente : le syst\`eme ne poss\`ede aucun \'etat
« en attente de validation » sur le chemin nominal.
\end{tcolorbox}

\section{Besoins non fonctionnels}
\begin{table}[H]
\centering
\caption{Exigences non fonctionnelles}
\label{tab:exigences-non-fonctionnelles}
\scriptsize
\begin{tabular}{@{}llp{8.4cm}@{}}
\toprule
\textbf{R\'ef.} & \textbf{Cat\'egorie} & \textbf{\'Enonc\'e et crit\`ere de v\'erification} \\
\midrule
ENF-01 & Souverainet\'e & Aucun appel sortant obligatoire ; fonctionnement nominal sans acc\`es Internet. \\
ENF-02 & Opposabilit\'e & Toute d\'ecision et toute action inscrites dans un journal dont la modification est techniquement emp\^ech\'ee. \\
ENF-03 & Tra\c cabilit\'e & Chaque d\'ecision porte ses sources documentaires, les r\`egles appliqu\'ees et les alternatives \'ecart\'ees. \\
ENF-04 & Performance & D\'elai entre r\'eception de l'\'ev\'enement et ex\'ecution de l'action inf\'erieur \`a [valeur] secondes au 95\textsuperscript{e} centile. \\
ENF-05 & Disponibilit\'e & Mode d\'egrad\'e : mise en file et rejeu int\'egral au retour du lien, sans perte. \\
ENF-06 & S\^uret\'e & Aucune action irr\'eversible ex\'ecutable en autonomie, quelle que soit la configuration. \\
ENF-07 & Portabilit\'e & D\'eploiement sur un serveur unique, sans service manag\'e externe. \\
ENF-08 & Maintenabilit\'e & Ajout d'une source de d\'etection sans modification du moteur d'orchestration. \\
ENF-09 & Utilisabilit\'e & Interface en fran\c cais ; la couleur ne porte jamais seule une information. \\
ENF-10 & S\'ecurit\'e & S\'eparation des r\^oles ; empreintes de mots de passe ; sessions \`a dur\'ee born\'ee. \\
\bottomrule
\end{tabular}
\end{table}

\section{Mod\`ele des cas d'utilisation}
\subsection{Diagramme de cas d'utilisation}
La figure~\ref{fig:cas-utilisation} pr\'esente le mod\`ele des cas d'utilisation.
Deux traits m\'eritent d'\^etre relev\'es avant lecture. D'une part, l'acteur
\emph{Source de d\'etection} et l'acteur \emph{Ordonnanceur} d\'eclenchent la cha\^ine
compl\`ete jusqu'\`a l'ex\'ecution : aucun cas d'utilisation humain ne s'ins\`ere entre la
d\'ecision et l'action. D'autre part, l'analyste n'appara\^it qu'en aval --- ses cas
d'utilisation sont tous post\'erieurs \`a une ex\'ecution.

\begin{figure}[H]
\centering
\begin{diagramme}[node distance=5mm]
  % ---- cas d'utilisation, colonne centrale
  \node (uc1) [cas=3.1cm] {Ing\'erer un \'ev\'enement normalis\'e};
  \node (uc2) [cas=3.1cm, below=6mm of uc1, umlaccentue] {Ex\'ecuter une action corrective en autonomie};
  \node (uc3) [cas=3.1cm, below=6mm of uc2] {\'Evaluer la r\'eversibilit\'e d'une action};
  \node (uc4) [cas=3.1cm, below=6mm of uc3, umlaccentue] {Annuler automatiquement une action};
  \node (uc5) [cas=3.1cm, below=6mm of uc4] {Notifier l'analyste a posteriori};
  \node (uc13)[cas=3.1cm, below=6mm of uc5, umlaccentue] {Confiner une menace hors catalogue};
  \node (uc6) [cas=3.1cm, below=6mm of uc13] {D\'eclencher une annulation manuelle};
  \node (uc7) [cas=3.1cm, below=6mm of uc6] {Consulter le journal d'audit};
  \node (uc8) [cas=3.1cm, below=6mm of uc7] {Consulter le portefeuille prioris\'e};
  \node (uc14)[cas=3.1cm, below=6mm of uc8] {Qualifier une menace et enrichir le catalogue};
  \node (uc9) [cas=3.1cm, below=6mm of uc14] {D\'efinir la politique de r\'eponse};
  \node (uc10)[cas=3.1cm, below=6mm of uc9] {G\'erer le catalogue de r\'eversibilit\'e};
  \node (uc11)[cas=3.1cm, below=6mm of uc10] {Configurer le coupe-circuit};
  \node (uc12)[cas=3.1cm, below=6mm of uc11] {Consulter la file du mode d\'egrad\'e};

  % ---- frontiere du systeme
  \begin{scope}[on background layer]
    \node (sys) [systeme, fit=(uc1)(uc12)] {};
    \node [font=\scriptsize\bfseries, anchor=north] at (sys.north) {CIRTDEFENSE};
  \end{scope}

  % ---- acteurs
  \umlacteur{src}{$(sys.west)+(-2.6cm,4.6cm)$}{Source de\\d\'etection}
  \umlacteur{ord}{$(sys.west)+(-2.6cm,1.6cm)$}{Ordonnanceur\\interne}
  \umlacteur{ana}{$(sys.east)+(2.6cm,1.2cm)$}{Analyste}
  \umlacteur{adm}{$(sys.east)+(2.6cm,-2.6cm)$}{Administrateur}
  \umlacteur{dec}{$(sys.west)+(-2.6cm,-3.4cm)$}{D\'ecideur}
  \umlacteur{eqp}{$(sys.east)+(2.6cm,4.6cm)$}{\'Equipement\\cible}

  % ---- associations acteur / cas
  \draw[association] (src) -- (uc1);
  \draw[association] (ord) -- (uc4);
  \draw[association] (ord) -- (uc12);
  \draw[association] (uc2) -- (eqp);
  \draw[association] (uc4) -- (eqp);
  \draw[association] (uc5) -- (ana);
  \draw[association] (ana) -- (uc6);
  \draw[association] (ana) -- (uc7);
  \draw[association] (ana) -- (uc8);
  \draw[association] (ana) -- (uc14);
  \draw[association] (adm) -- (uc9);
  \draw[association] (adm) -- (uc10);
  \draw[association] (adm) -- (uc11);
  \draw[association] (adm) -- (uc12);
  \draw[association] (dec) -- (uc7);
  \draw[association] (dec) -- (uc8);

  % ---- relations entre cas
  \draw[inclut] (uc1) to[out=0,in=90] node[role, pos=0.5]{\guillemotleft include\guillemotright} (uc2.north east);
  \draw[inclut] (uc2) -- node[role]{\guillemotleft include\guillemotright} (uc3);
  \draw[inclut] (uc4) -- node[role]{\guillemotleft include\guillemotright} (uc3);
  \draw[inclut] (uc2) to[out=180,in=180] node[role,left=1mm]{\guillemotleft include\guillemotright} (uc5);
  \draw[etend]  (uc13) to[out=180,in=180] node[role,left=1mm]{\guillemotleft extend\guillemotright} (uc2);
  \draw[etend]  (uc14) to[out=0,in=0] node[role,right=1mm]{\guillemotleft extend\guillemotright} (uc13);
\end{diagramme}
\caption{Diagramme de cas d'utilisation du syst\`eme}
\label{fig:cas-utilisation}
\end{figure}

\subsection{Description d\'etaill\'ee des cas critiques}
Les fiches qui suivent adoptent le format usuel : identification, acteurs,
pr\'econditions, sc\'enario nominal, sc\'enarios alternatifs, postconditions, exigences
couvertes. Seuls les cas portant l'apport du travail sont d\'etaill\'es ici ; les
autres figurent en annexe~\ref{ann:cas-utilisation}.

\begin{fichecas}{UC-02}{Ex\'ecuter une action corrective en autonomie}
\rubrique{Acteur principal}{Source de d\'etection (acteur syst\`eme).}
\rubrique{Acteurs secondaires}{\'Equipement cible, Analyste (notifi\'e apr\`es coup).}
\rubrique{Objectif}{Contenir une menace sans attendre d'intervention humaine.}
\rubrique{Pr\'econditions}{Un \'ev\'enement normalis\'e a \'et\'e ing\'er\'e et rattach\'e \`a un
  incident ; le coupe-circuit est ferm\'e ; le syst\`eme n'est pas en mode d\'egrad\'e.}
\rubrique{Sc\'enario nominal}{%
  \begin{enumerate}
    \item Le syst\`eme enrichit le contexte de l'incident \`a partir du fonds documentaire.
    \item Il v\'erifie que le contexte est document\'e \emph{(EF-04)}.
    \item Il classe l'incident au catalogue m\'etier et s\'electionne le playbook associ\'e.
    \item Il param\`etre les actions \`a partir des indicateurs de l'\'ev\'enement.
    \item Il \'ecarte toute action non r\'eversible \emph{(EF-14, cf. UC-03)}.
    \item Il soumet les actions restantes \`a la politique compil\'ee \emph{(EF-15)}.
    \item Il ex\'ecute les actions autoris\'ees aupr\`es des actionneurs \emph{(EF-07)}.
    \item Il inscrit d\'ecision et r\'esultats au journal d'audit.
    \item Il notifie l'analyste \emph{(EF-13, cf. UC-05)}.
  \end{enumerate}}
\rubrique{A1 --- contexte non document\'e}{\`A l'\'etape 2, le syst\`eme n'ex\'ecute rien,
  inscrit un refus motiv\'e et \'emet une alerte indiquant que la menace n'est pas
  contenue. Reprise possible par UC-13.}
\rubrique{A2 --- toutes les actions refus\'ees}{\`A l'\'etape 6, aucune action n'est
  autoris\'ee ; le syst\`eme inscrit le verdict de politique et notifie l'analyste.}
\rubrique{A3 --- coupe-circuit ouvert}{Le cas se termine sans ex\'ecution ; l'incident
  reste ouvert et remonte dans le portefeuille.}
\rubrique{E1 --- \'echec d'un actionneur}{L'action est marqu\'ee en \'echec ; les autres
  se poursuivent ; l'\'echec alimente le compteur du coupe-circuit.}
\rubrique{Postconditions}{Les actions ex\'ecut\'ees portent chacune un jeton
  d'annulation ; l'incident passe \`a l'\'etat \emph{confin\'e} ; la boucle de contr\^ole
  post-action est arm\'ee sur la cible.}
\rubrique{Exigences}{EF-05, EF-06, EF-07, EF-13, EF-14, EF-15, EF-26.}
\end{fichecas}

\begin{fichecas}{UC-04}{Annuler automatiquement une action}
\rubrique{Acteur principal}{Ordonnanceur interne.}
\rubrique{Objectif}{Rattraper une action correcte du point de vue de la menace mais
  nuisible du point de vue de la cible.}
\rubrique{Pr\'econditions}{Au moins une action a \'et\'e ex\'ecut\'ee et porte un jeton
  d'annulation ; une mesure de r\'ef\'erence de la cible a \'et\'e prise avant l'action.}
\rubrique{Sc\'enario nominal}{%
  \begin{enumerate}
    \item L'ordonnanceur d\'eclenche la boucle de contr\^ole.
    \item Le syst\`eme mesure l'\'etat courant de chaque cible ayant subi une action.
    \item Il compare cet \'etat \`a la mesure de r\'ef\'erence.
    \item Si la d\'egradation d\'epasse le seuil, il d\'eclenche l'annulation au moyen du jeton.
    \item Il inscrit l'annulation au journal et notifie l'analyste.
  \end{enumerate}}
\rubrique{A1 --- absence de r\'ef\'erence}{Le syst\`eme s'abstient : sans mesure
  ant\'erieure, la comparaison n'aurait aucun sens. L'abstention est motiv\'ee et inscrite.}
\rubrique{A2 --- am\'elioration ou stabilit\'e}{Aucune annulation ; l'action est confirm\'ee.}
\rubrique{E1 --- \'echec de l'annulation}{L'action passe \`a \emph{\'echec d'annulation} ;
  une alerte de gravit\'e critique est \'emise : c'est le seul cas o\`u une intervention
  humaine imm\'ediate est requise.}
\rubrique{Postconditions}{L'\'etat de la cible est restaur\'e ou l'incapacit\'e \`a le
  restaurer est trac\'ee et signal\'ee.}
\rubrique{Exigences}{EF-25, EF-21, EF-22.}
\end{fichecas}

\begin{fichecas}{UC-13}{Confiner une menace hors catalogue}
\rubrique{Acteur principal}{Source de d\'etection.}
\rubrique{Objectif}{\'Eviter le silence face \`a une menace qu'aucun playbook ne couvre.}
\rubrique{Pr\'econditions}{Le contexte de l'incident n'est pas document\'e : UC-02 s'est
  arr\^et\'e au sc\'enario alternatif A1.}
\rubrique{Sc\'enario nominal}{%
  \begin{enumerate}
    \item Le syst\`eme extrait les indicateurs observables de l'\'ev\'enement.
    \item Pour chaque indicateur exploitable, il d\'eduit un geste de confinement.
    \item Il partage les gestes en deux ensembles selon leur r\'eversibilit\'e.
    \item Il ex\'ecute sans attendre les gestes r\'eversibles \`a rayon d'impact restreint.
    \item Il inscrit les gestes \`a effet durable comme \emph{en attente de d\'ecision humaine}.
    \item Il \'emet une alerte persistante d\'ecrivant ce qui a \'et\'e fait et ce qui reste \`a d\'ecider.
  \end{enumerate}}
\rubrique{R\`egle de gestion}{Le syst\`eme ne d\'eduit jamais un geste d'un \emph{type
  d'attaque suppos\'e} mais d'un \emph{indicateur constat\'e}. Cette distinction est ce
  qui permet de sortir du catalogue sans lever l'exigence de non-invention.}
\rubrique{A1 --- aucun indicateur exploitable}{Le syst\`eme n'ex\'ecute rien et \'emet une
  alerte d'abstention motiv\'ee.}
\rubrique{A2 --- adresse interne}{Une adresse appartenant au plan d'adressage interne
  n'est jamais bloqu\'ee : le geste serait plus dommageable que la menace.}
\rubrique{Postconditions}{L'incident est ouvert, partiellement confin\'e, et porte une
  demande de qualification (UC-14).}
\rubrique{Exigences}{EF-27, EF-28, EF-14, EF-15.}
\end{fichecas}

\section{Conclusion du chapitre}
Les besoins recueillis \'etablissent une contrainte dominante : l'autonomie n'est
acceptable que si elle est \emph{born\'ee}, et la borne retenue est la
r\'eversibilit\'e. Le chapitre suivant traduit cette contrainte en structures :
c'est elle qui explique la plupart des choix de conception, jusqu'\`a la forme des
classes du domaine.
